\section{How to run it yourself}
\label{sec:running}

These commands are cross-checked against \file{docs/REPRODUCING.md} and the
\file{Makefile}. Run them from the repository root.

\subsection{One-time setup}

\begin{lstlisting}[style=shellstyle,caption={Create an isolated environment and install the package.}]
# 1. Create an isolated Python environment (so you don't clutter your system)
python3 -m venv .venv
source .venv/bin/activate        # Windows: .venv\Scripts\activate

# 2. Install the package + dev tools (pytest, jupyterlab, matplotlib, ...)
make install                     # same as: pip install -e ".[dev]"
\end{lstlisting}

A \textbf{virtual environment} (\code{venv}) is just a private sandbox of Python
packages for this project. \code{pip install -e .} installs the \code{lunar}
package in ``editable'' mode, meaning Python sees your local \file{lunar/} folder
directly --- edit a file and the change takes effect immediately.

\subsection{Run the tests first (always do this)}

\begin{lstlisting}[style=shellstyle]
python3 -m pytest -q
\end{lstlisting}

This runs the suite in \file{tests/} (grid, properties, solver, equilibrium,
bedrock, ephem, constants). They finish in under a minute and should
\textbf{all pass}. Reading a test is often the clearest example of how a function
is meant to be used.

\subsection{Verify the committed headline numbers (no computation needed)}

The canonical result ships in the repo, so you can check it instantly:

\begin{lstlisting}[caption={Verify the committed headline \Kdstar\ values.}]
import json, math
r = json.load(open('output/kd_retrieval_results.json'))
print('A15 K_d* =', r['A15']['kd_star'] * 1e3, 'mW/m/K')
print('A17 K_d* =', r['A17']['kd_star'] * 1e3, 'mW/m/K')
assert math.isclose(r['A15']['kd_star']*1e3, 4.58, abs_tol=0.01)
assert math.isclose(r['A17']['kd_star']*1e3, 8.12, abs_tol=0.01)
print('Headline values verified.')
\end{lstlisting}

\subsection{Reproduce the science (recompute from scratch)}

\begin{lstlisting}[style=shellstyle,caption={The one-word \file{Makefile} targets.}]
make help          # list every one-word command
make retrieve      # the core K_d retrieval + bootstrap (~5 min)
                   #   -> writes output/kd_retrieval_results.json
make aux           # all auxiliary sweeps, model selection, error budget, MCMC
make figures       # regenerate every figure
make paper         # compile the PDFs (needs a LaTeX install)
make all           # everything end-to-end
\end{lstlisting}

\code{make retrieve} is the one to run first --- it regenerates the headline JSON.
The slow auxiliary sweeps (\code{make aux}) are rarely needed because their
results are already committed.

\subsection{Run a single illuminating script}

To \emph{see} the deep-temperature physics (\autoref{sec:skin}) with your own
eyes:

\begin{lstlisting}[style=shellstyle]
python3 scripts/analysis/phase2_depth_convergence.py
\end{lstlisting}

It prints the skin depth, the deep gradient, and confirms the mean temperature
keeps rising along $Q_b/K$, then saves
\file{output/figures/fig\_phase2\_depth\_convergence.png}
(\autoref{fig:phase2}). Open that image. This is one of the best on-ramps to
understanding the project.

\begin{tryit}
Regenerate the five teaching figures in this very handbook. From the repository
root:
\begin{lstlisting}[style=shellstyle]
export MPLCONFIGDIR=/tmp/mpl
python3 docs/teaching/make_teaching_figures.py
\end{lstlisting}
This re-runs the certified equilibrium solver for both sites and re-reads the
committed results JSON, writing fresh PNG$+$PDF assets into
\file{docs/teaching/figures/}. Then rebuild this PDF with
\code{latexmk -pdf beginners\_guide.tex}.
\end{tryit}

\subsection{The notebooks (guided, in order)}

\begin{lstlisting}[style=shellstyle]
jupyter lab notebooks/
\end{lstlisting}

Run them in order: \code{00\_setup} (checks), \code{01\_methods} (figures $+$
parameter table), \code{02\_retrieval} (the \Kd{} retrieval $+$ bootstrap),
\code{03\_results}, \code{04\_discussion}, \code{05\_animations}. There's also
\code{equilibrium\_demo.ipynb}, a standalone demo that a brute-force
1000-lunation spin-up converges onto the fast Anchor Point solver --- a great way
to \emph{believe} \autoref{sec:anchor}.

\subsection{Where outputs land}

\begin{itemize}[leftmargin=1.4em,itemsep=2pt]
  \item \file{output/*.json} --- numerical results (the canonical
    \file{kd\_retrieval\_results.json} and the auxiliary JSONs).
  \item \file{output/figures/*.png} / \file{*.pdf} --- diagnostic and appendix
    figures.
  \item \file{paper/letter/figures/*.pdf} --- the paper's figures.
  \item \file{paper/letter/letter.pdf} --- the compiled manuscript.
\end{itemize}
